\documentclass[../../main.tex]{subfiles}
\begin{document}

{\bf [解答]}

\paragraph{(1) 正しい}

対偶をとると「$a > b \implies a \neq b$」だからこれは明らかである．

\paragraph{(2) 正しくない}

反例を構成する．座標平面で $A(0,0)$, $B(1,0)$, $C(1,1)$, $D(0,1)$ とすると, どの3点も一直線上にないが,
\begin{align}
  \overrightarrow{AB} + \overrightarrow{CD} =
  \begin{pmatrix} 1 \\ 0
  \end{pmatrix} +
  \begin{pmatrix} -1 \\ 0
  \end{pmatrix} = \vec{0}
\end{align}
である.

\paragraph{(3) 正しい}

まず$A=B$を示す. 題意より$A \subset B$ である. 次に,題意より
\begin{align}
  B \subset C \subset A \implies B \subset A
\end{align}
である．よって $A=B$ が従う．同様にして $B=C$ も示されるので
\begin{align}
  A=B=C
\end{align}
である．

\paragraph{(4) 正しい}

数学的帰納法で示す. $n=5$の時は$n^2=25, 2^n=32$ より $25 < 32$となって成立する.
そこで$n=k \, (k \in \mathbb{N}, k \geqq 5)$ での成立を仮定する．
$n=k+1$の時はまず
\begin{align}
  2^{k+1} > 2 \cdot k^2 \quad \text{(仮定)}
\end{align}
が成り立つ．この右辺をさらに下から評価すると
\begin{align}
  2k^2 > (k+1)^2
  \iff
  k^2-2k-1>0
\end{align}
だが，これは$k \ge 5$より成立する．以上より
\begin{align}
  2^{k+1} > (k+1)^2
\end{align}
となって $n=k+1$ でも題意は成立する. 以上から示された.

\paragraph{(5) 正しくない}

背理法により題意を示す．
題意が正しいとすると $\dfrac{1}{N} = 0$ なる $N \in \mathbb{N}$ が存在する. 両辺に $N (\neq 0)$ をかけると
\begin{align}
  1 = 0
\end{align}
となって矛盾するから示された.

\newpage

\end{document}
